\section{Introdução}
\label{sec:introducao}

A transformação digital no setor público brasileiro consolidou a oferta de serviços governamentais por meio de arquiteturas multicanal, permitindo que cidadãos acessem informações críticas tanto por portais web tradicionais quanto por aplicativos móveis nativos ou híbridos. No âmbito do Distrito Federal, o Departamento de Trânsito (Detran-DF) disponibiliza a consulta ao histórico de infrações de trânsito por meio do Portal de Serviços (``Detran Digital'') \cite{detrandf2026portal} e do aplicativo móvel ``DETRAN-DF'' \cite{detrandf2026app}, atendendo a uma demanda contínua de condutores e proprietários de veículos.

Embora ambos os canais compartilhem o mecanismo de autenticação e integrem dados locais e do Registro Nacional de Infrações de Trânsito (RENAINF), eles operam sobre clientes tecnológicos distintos e consomem rotas de serviços de \textit{backend} específicas. Essa separação arquitetural introduz o problema central investigado neste estudo: a possibilidade de assimetria de qualidade entre os canais, manifestada por divergências nos campos exibidos de um mesmo auto de infração, diferenças no esforço operacional e no consumo de dados móveis, ou falhas de conclusão da consulta sob diferentes dispositivos e condições de conectividade. Tal inconsistência impacta diretamente a tomada de decisão do cidadão quanto ao pagamento de multas, adesão a descontos ou interposição de recursos dentro dos prazos legais estipulados pelo Código de Trânsito Brasileiro \cite{brasil1997ctb}.

O objetivo deste trabalho consiste em definir os requisitos e especificar formalmente o plano de avaliação comparativa da qualidade do produto de software para o serviço de consulta de infrações do Detran-DF (Portal de Serviços $\times$ Aplicativo), executando as Fases 1 e 2 do processo de avaliação preconizado pela norma ISO/IEC 25040 \cite{iso25040_2011}. A avaliação adota uma perspectiva estritamente externa (caixa-preta), voltada a apoiar decisões técnicas e gerenciais de priorização de melhorias e recomendação de uso por contexto.

As principais contribuições desta entrega compreendem: (i) o mapeamento de partes interessadas, decisões institucionais e delimitação de escopo em caixa-preta alinhados aos Objetivos de Desenvolvimento Sustentável (ODS 16 e ODS 9) \cite{un2015agenda}; (ii) a adaptação fundamentada dos modelos da família SQuaRE, combinando Adequação Funcional e Eficiência de Desempenho da ISO/IEC 25010:2023 \cite{iso25010_2023} com a Cobertura de Contexto da ISO/IEC 25010:2011 \cite{iso25010_2011} e ISO/IEC 25022:2016 \cite{iso25022_2016}; e (iii) a derivação completa de uma árvore \textit{Goal Question Metric} (GQM) \cite{basili1994gqm} contendo objetivos, questões, hipóteses, métricas reprodutíveis e regras de julgamento, acompanhada de governança de projeto baseada no \textit{Practical Software Measurement} (PSM/CID) \cite{mcgarry2002psm}.

O restante deste artigo está organizado da seguinte forma: a Seção~\ref{sec:referencial-teorico} fundamenta o referencial teórico unificado sobre a família SQuaRE, a abordagem GQM e o modelo gerencial PSM/CID; a Seção~\ref{sec:metodologia} apresenta a visão geral da metodologia de avaliação e detalha passo a passo a Fase 1 (Definição dos Requisitos de Avaliação) e a Fase 2 (Especificação da Avaliação); por fim, a Seção~\ref{sec:declaracao-ia} registra a declaração de uso ético e validado de ferramentas de inteligência artificial, seguida do apêndice de contribuição individual da equipe.